% =============================================================================
% 4.2 DIAGRAMAS DE CLASES POR DOMINIOS CRÍTICOS
% =============================================================================
\section{Diagramas de Clases por Dominios Críticos}
\label{sec:clases}

Esta sección documenta la \textbf{estructura estática} del backend: cómo se organizan las
clases dentro de cada dominio y cómo fluyen los datos entre ellas. El patrón es uniforme
y deliberado, lo que reduce la carga cognitiva al navegar el código: \textit{todo dominio
se lee de la misma forma}.

\subsection{El patrón de capas por \textit{feature}}
\label{subsec:patron_capas}

Cada dominio implementa tres responsabilidades bien separadas, según el principio de
responsabilidad única (SRP):

\begin{itemize}[noitemsep]
    \item \comando{@RestController}: frontera HTTP. Recibe DTOs, valida y delega; nunca
    contiene lógica de negocio ni acceso a datos.
    \item \comando{@Service}: lógica de negocio y límite transaccional
    (\comando{@Transactional}). Orquesta repositorios y servicios externos.
    \item \comando{Repository}: acceso a datos tipado con jOOQ a través de un único
    \comando{DSLContext} inyectado por Spring.
\end{itemize}

La figura~\ref{fig:clases-capas} ilustra el patrón sobre tres dominios críticos
—\texttt{profile}, \texttt{portfolio} y \texttt{cvstudio}— y evidencia la dependencia
descendente Controller~$\rightarrow$~Service~$\rightarrow$~Repository.

El backend organiza el código en \textbf{paquetes por dominio} bajo
\pathbreak{com.ethoshub}, cada uno autocontenido. La figura~\ref{fig:paquetes-backend}
ofrece un mapa de los paquetes reales, agrupados por su rol funcional.

\vspace{0.3cm}
\begin{figure}[h!]
\centering
\begin{tikzpicture}[
    pkg/.style={rectangle, rounded corners=2pt, draw=c4Sistema!70!black, fill=c4Componente!28,
        font=\sffamily\scriptsize, align=center, minimum width=2.5cm, minimum height=0.7cm, inner sep=2pt},
    grp/.style={rectangle, rounded corners=4pt, draw=grisISO!60, dashed, inner sep=7pt},
    lbl/.style={font=\sffamily\scriptsize\bfseries, text=colorPrimario}]
    % Identidad
    \node[pkg] (auth) at (0,3) {auth};
    \node[pkg] (sec) at (2.6,3) {security};
    \node[pkg] (prof) at (5.2,3) {profile};
    \node[pkg] (pref) at (7.8,3) {preferences};
    \node[lbl] at (3.9,3.7) {Identidad y cuentas};
    \begin{scope}[on background layer]
      \node[grp, fit=(auth)(pref), draw=c4Persona!60] {};
    \end{scope}
    % Contenido
    \node[pkg] (proj) at (0,1.4) {project};
    \node[pkg] (skill) at (2.6,1.4) {skills};
    \node[pkg] (exp) at (5.2,1.4) {experience};
    \node[pkg] (edu) at (7.8,1.4) {education};
    \node[pkg] (port) at (1.3,0.4) {portfolio};
    \node[pkg] (cv) at (3.9,0.4) {cvstudio};
    \node[lbl] at (3.9,2.1) {Contenido profesional};
    % Interacciones
    \node[pkg] (chat) at (0,-1.1) {chat};
    \node[pkg] (notif) at (2.6,-1.1) {notifications};
    \node[pkg] (conn) at (5.2,-1.1) {connections};
    \node[pkg] (recr) at (7.8,-1.1) {recruiter (+ ai)};
    \node[lbl] at (3.9,-0.4) {Interacciones};
    % Transversal
    \node[pkg, fill=colorAcento!20] (shared) at (1.3,-2.4) {shared};
    \node[pkg, fill=colorAcento!20] (config) at (3.9,-2.4) {config};
    \node[pkg, fill=colorAcento!20] (admin) at (6.5,-2.4) {admin};
    \node[lbl] at (3.9,-1.8) {Transversal};
\end{tikzpicture}
\caption{Mapa de paquetes del backend bajo \texttt{com.ethoshub}, agrupados por rol.}
\label{fig:paquetes-backend}
\end{figure}

\vspace{0.3cm}
\begin{figure}[h!]
\centering
\begin{tikzpicture}[
    cls/.style={rectangle, draw=c4Sistema!70!black, fill=c4Componente!25, font=\sffamily\scriptsize,
        align=center, minimum width=2.7cm, minimum height=0.8cm, inner sep=3pt},
    ctrl/.style={cls, fill=c4Componente!50},
    rel/.style={-{Stealth[length=2mm]}, semithick, draw=grisISO},
    node distance=0.7cm and 0.5cm
]
    \node[ctrl] (pc) at (0,3) {\texttt{ProfileController}};
    \node[cls] (ps) at (0,1.8) {\texttt{ProfileService}};
    \node[cls] (pr) at (0,0.6) {\texttt{ProfileRepository}};
    \node[ctrl] (bc) at (3.6,3) {\texttt{PortfolioSettings}\\\texttt{Controller}};
    \node[cls] (bs) at (3.6,1.8) {\texttt{PortfolioSettings}\\\texttt{Service}};
    \node[cls] (br) at (3.6,0.6) {\texttt{PortfolioRepository}};
    \node[ctrl] (cc) at (7.2,3) {\texttt{CvDocument}\\\texttt{Controller}};
    \node[cls] (cs) at (7.2,1.8) {\texttt{CvDocumentService}\\\texttt{+ LatexCompileSvc}};
    \node[cls] (cr) at (7.2,0.6) {\texttt{CvDocumentRepository}};

    \node[cls, fill=c4Datos!25, draw=c4Datos!70!black, minimum width=8.5cm] (dsl) at (3.6,-0.7)
        {jOOQ \texttt{DSLContext} $\rightarrow$ PostgreSQL \texttt{core}};

    \foreach \a/\b in {pc/ps, ps/pr, bc/bs, bs/br, cc/cs, cs/cr} \draw[rel] (\a) -- (\b);
    \draw[rel] (pr) -- (dsl); \draw[rel] (br) -- (dsl); \draw[rel] (cr) -- (dsl);
\end{tikzpicture}
\caption{Patrón de capas por \textit{feature} en los dominios \texttt{profile}, \texttt{portfolio} y \texttt{cvstudio}.}
\label{fig:clases-capas}
\end{figure}

El dominio \texttt{cvstudio} es el de mayor riqueza estructural: su servicio coordina no
solo el repositorio, sino dos colaboradores especializados —\comando{GeminiAiService}
(asistencia de IA) y \comando{LatexCompileService} (compilación a PDF)—, ambos
documentados en la sección~\ref{sec:secuencias} (pág.~\pageref{sec:secuencias}).

\subsection{Modelado de DTOs y registros tipados de jOOQ}
\label{subsec:dtos}

Los \textbf{objetos de transferencia de datos} (DTOs) se implementan como
\comando{record} de Java inmutables, con validación declarativa mediante \textbf{Jakarta
Bean Validation}. Se distinguen dos familias:

\begin{itemize}[noitemsep]
    \item \textbf{DTOs de petición} (\texttt{*Request}): anotados con restricciones
    (\comando{@NotBlank}, \comando{@Email}, \comando{@Size}) que Spring evalúa
    automáticamente al deserializar el cuerpo JSON.
    \item \textbf{DTOs de respuesta} (\texttt{*Response} / \texttt{*Dto}): proyecciones de
    solo lectura que el servicio construye a partir de los registros tipados de jOOQ.
\end{itemize}

Entre ambos median los \textbf{registros generados por jOOQ} a partir del esquema real de
la base de datos: clases como \texttt{ProfilesBasicRecord} o \texttt{CvDocumentsRecord}
que ofrecen acceso tipado y verificado en compilación a cada columna, eliminando los
errores de SQL por cadenas mágicas.

La figura~\ref{fig:dto-flujo} muestra el viaje de los datos desde el DTO de petición
validado hasta el DTO de respuesta serializado.

\vspace{0.3cm}
\begin{figure}[h!]
\centering
\begin{tikzpicture}[node distance=0.6cm and 1.3cm,
    box/.style={rectangle, rounded corners=2pt, draw, font=\sffamily\scriptsize, align=center,
        minimum width=2.6cm, minimum height=0.85cm, inner sep=3pt},
    rel/.style={-{Stealth[length=2mm]}, semithick}]
    \node[box, fill=c4Componente!40] (req) {\texttt{RegisterRequest}\\\scriptsize @NotBlank @Email @Size};
    \node[box, fill=colorAcento!30, right=of req] (val) {Bean Validation\\\scriptsize (Spring)};
    \node[box, fill=c4Componente!40, right=of val] (svc) {\texttt{AuthService}\\\scriptsize lógica};
    \node[box, fill=c4Datos!25, below=of svc] (db) {jOOQ\\\texttt{ProfilesBasicRecord}};
    \node[box, fill=c4Componente!40, below=of req] (resp) {\texttt{AuthResponse}\\\scriptsize token + rol};

    \draw[rel] (req) -- node[c4etiqueta, above]{JSON} (val);
    \draw[rel] (val) -- node[c4etiqueta, above]{válido} (svc);
    \draw[rel] (svc) -- (db);
    \draw[rel] (svc) to[bend left=10] (resp);
    \draw[rel, dashed] (val) to[bend right=30] node[c4etiqueta, left]{400 +\\errores} (resp);
\end{tikzpicture}
\caption{Flujo de un DTO de petición validado hacia el DTO de respuesta.}
\label{fig:dto-flujo}
\end{figure}

\begin{lstlisting}[language=Java, caption={DTO de registro con validación declarativa (\texttt{auth/dto/RegisterRequest.java}).}, label={lst:registerrequest}]
public record RegisterRequest(
    @NotBlank @Email                 String email,
    @NotBlank @Size(min = 8)         String password,
    @NotNull                         UserRole role,   // PROFESSIONAL | RECRUITER
    String firstName, String lastName,
    String phoneCode, String phoneNumber
) { }
\end{lstlisting}

Cuando la validación falla, \comando{MethodArgumentNotValidException} es interceptada por
el manejador global (sección~\ref{sec:excepciones}, pág.~\pageref{sec:excepciones}) y
traducida a una respuesta \texttt{400} con la lista de errores por campo, sin que el
controlador escriba una sola línea de código de validación manual.

\begin{NotaTecnica}
La elección de \texttt{record} inmutables para los DTOs no es estética: garantiza que un
objeto de petición no pueda mutar tras su validación, eliminando una clase entera de
errores (\textit{time-of-check to time-of-use}). La inmutabilidad también facilita el
razonamiento concurrente, relevante para el modelo \textit{stateless} del servidor.
\end{NotaTecnica}

\subsection{Resumen de DTOs por dominio crítico}
\label{subsec:dtos_resumen}

La tabla~\ref{tab:dtos} cataloga los DTOs más representativos y su rol en el contrato de
la API.

\vspace{0.2cm}
\noindent
\renewcommand{\arraystretch}{1.3}
\fontsize{8.5}{10.2}\selectfont\sffamily
\begin{tabularx}{\textwidth}{|L{4.2cm}|C{1.8cm}|Y|}
\hline
\rowcolor{headerTabla}
\sffamily\textbf{DTO} & \sffamily\textbf{Familia} & \sffamily\textbf{Uso} \\
\hline
\texttt{RegisterRequest} & Petición & Alta de usuario con rol y datos de contacto. \\ \hline
\rowcolor{grisTecnico}
\texttt{LoginRequest} & Petición & Credenciales para \textit{login} con email/contraseña. \\ \hline
\texttt{AuthResponse} & Respuesta & Token, tipo de token, expiración y rol resuelto. \\ \hline
\rowcolor{grisTecnico}
\texttt{CvDocumentRequest} & Petición & Título y código LaTeX del CV a crear/actualizar. \\ \hline
\texttt{AiAssistRequest} & Petición & Prompt, contenido, modo e historial para Gemini. \\ \hline
\rowcolor{grisTecnico}
\texttt{TalentDto} & Respuesta & Proyección de un perfil profesional para reclutadores. \\ \hline
\texttt{ChatMessageRequest} & Petición & Mensaje, \texttt{chatId} y clave de idempotencia. \\ \hline
\end{tabularx}
\label{tab:dtos}

\subsection{Diagrama de clases UML del dominio \texttt{cvstudio}}
\label{subsec:uml_cvstudio}

La figura~\ref{fig:uml-cvstudio} ofrece una vista UML del dominio más rico del backend,
con sus atributos y operaciones públicas. Evidencia cómo \comando{CvDocumentService}
coordina tres colaboradores y traduce entre la entidad de dominio y los DTOs.

\clearpage
\vspace{0.3cm}
\begin{figure}[h!]
\centering
\begin{tikzpicture}[
    uml/.style={rectangle, draw=c4Sistema!70!black, fill=white, font=\sffamily\scriptsize,
        align=left, inner sep=4pt, minimum width=3.6cm},
    titulo/.style={font=\sffamily\scriptsize\bfseries, text=colorPrimario},
    rel/.style={-{Stealth[length=2mm]}, semithick, draw=grisISO},
    agg/.style={-{Diamond[open, length=3mm]}, semithick, draw=grisISO}]

    \node[uml] (ctrl) at (0,3.4) {%
        {\titulo CvDocumentController}\\\hrule height0.4pt\vspace{2pt}
        + list() : ApiResponse\\
        + create(req) : ApiResponse\\
        + aiAssist(req) : ApiResponse\\
        + compileLatex(req) : byte[]};
    \node[uml] (svc) at (0,0.6) {%
        {\titulo CvDocumentService}\\\hrule height0.4pt\vspace{2pt}
        - repo : CvDocumentRepository\\
        - gemini : GeminiAiService\\
        - latex : LatexCompileService\\\hrule height0.4pt\vspace{2pt}
        + assist(...) : String\\
        + compile(...) : byte[]};
    \node[uml, fill=c4Datos!12] (repo) at (-3.6,-2.4) {%
        {\titulo CvDocumentRepository}\\\hrule height0.4pt\vspace{2pt}
        - dsl : DSLContext\\
        + findByProfile(id)\\
        + upsert(rec)};
    \node[uml, fill=colorAcento!10] (gem) at (0,-2.4) {%
        {\titulo GeminiAiService}\\\hrule height0.4pt\vspace{2pt}
        + generate\\Assistance(...)};
    \node[uml, fill=colorAcento!10] (tex) at (3.6,-2.4) {%
        {\titulo LatexCompileService}\\\hrule height0.4pt\vspace{2pt}
        + compile(code,\\\ \ fotoUrl)};

    \draw[rel] (ctrl) -- (svc);
    \draw[agg] (svc) -- (repo);
    \draw[agg] (svc) -- (gem);
    \draw[agg] (svc) -- (tex);
\end{tikzpicture}
\caption{Diagrama de clases UML del dominio \texttt{cvstudio} (atributos y operaciones).}
\label{fig:uml-cvstudio}
\end{figure}

\subsection{Acceso a datos tipado con jOOQ}
\label{subsec:jooq}

Los repositorios no escriben SQL como cadenas de texto: usan el \textbf{DSL tipado} de
jOOQ, generado a partir del esquema real de la base de datos. El \textit{listado}
\ref{lst:jooq} muestra un repositorio representativo donde cada nombre de tabla y columna
está verificado en tiempo de compilación.

\begin{lstlisting}[language=Java, caption={Acceso tipado con jOOQ en un repositorio (patrón real).}, label={lst:jooq}]
@Repository
public class CvDocumentRepository {
    private final DSLContext dsl;
    public List<CvDocumentsRecord> findByProfile(UUID profileId) {
        return dsl.selectFrom(CV_DOCUMENTS)
                  .where(CV_DOCUMENTS.PROFILE_ID.eq(profileId))
                  .and(CV_DOCUMENTS.DELETED_AT.isNull())   // borrado logico
                  .orderBy(CV_DOCUMENTS.UPDATED_AT.desc())
                  .fetch();
    }
}
\end{lstlisting}

\clearpage
\begin{NotaTecnica}
La ventaja decisiva de jOOQ sobre el SQL textual o un ORM por reflexión es la
\textbf{seguridad en compilación}: si una columna se renombra en la base de datos y se
regeneran las clases, cualquier consulta que la referencie deja de compilar. Esto convierte
una clase entera de errores de tiempo de ejecución en errores detectados antes del
despliegue, en línea con la estrategia de verificación del capítulo~8.
\end{NotaTecnica}
